
\section{Methods}

\subsection{Data Scope}
We work with aggregated data only. For each split we use selection counts and rates, per-group confusion counts, and calibration-bin summaries. No row-level data is required for the analyses in this whitepaper.

\subsection{Estimators and Uncertainty}
Selection-rate proportions use Wilson 95\% confidence intervals \citep{brown2001interval}. For the disparity ratio (\cref{eq:air}), we use a deterministic delta-method confidence interval on $\log(\mathrm{AIR})$; approval-rate gaps (\cref{eq:srg}) use the conservative difference of separate Wilson interval endpoints, $[L_g-U_{\mathrm{ref}},\,U_g-L_{\mathrm{ref}}]$. That SRG endpoint range is not itself a calibrated 95\% confidence interval. Two-sided p-values use a two-proportion z-test by default, switching to Fisher's exact test under a small-sample policy (small group size or sparse $2\times 2$ cells).

Historical and amplification intervals and p-values are conditional on this single generated fixture and its configured row count; they do not include fixture-generator or across-run uncertainty. The policy-determined intrinsic interval and p-value fields are non-inferential and are retained only for format symmetry.

\subsection{Multiplicity Adjustment}
For multi-class race, we perform pairwise comparisons under the recorded effective-reference policy. P-values across the race-pair family are adjusted using the Holm--Bonferroni step-down procedure at level $\alpha=0.05$. Confidence intervals are pairwise and are not multiplicity-adjusted; therefore their apparent decisions can differ from the adjusted p-values.
